% ============================================================
% Seminário - Software Básico
% Grupo 2: Ligador (Linker), Bibliotecas e Ligação Estática
% ============================================================
\documentclass[aspectratio=169, 11pt]{beamer}

% ---- Tema Moloch ----
\usetheme[progressbar=frametitle, numbering=fraction, block=fill]{moloch}

% ---- Paleta acadêmica ----
\definecolor{primary}{RGB}{20,40,72}        % azul marinho institucional
\definecolor{accent}{RGB}{0,105,148}        % azul petróleo
\definecolor{lightbg}{RGB}{245,246,248}     % cinza claro para blocos
\definecolor{codebg}{RGB}{240,241,244}      % fundo de código
\definecolor{success}{RGB}{40,110,50}       % verde
\definecolor{danger}{RGB}{170,40,40}        % vermelho

\setbeamercolor{normal text}{fg=primary, bg=white}
\setbeamercolor{alerted text}{fg=accent}
\setbeamercolor{frametitle}{bg=primary, fg=white}
\setbeamercolor{progress bar}{fg=accent, bg=primary!20}
\setbeamercolor{block title}{fg=white, bg=primary!85}
\setbeamercolor{block body}{fg=primary, bg=lightbg}
\setbeamercolor{block title alerted}{fg=white, bg=danger}
\setbeamercolor{block body alerted}{fg=primary, bg=danger!5}
\setbeamercolor{block title example}{fg=white, bg=success!85}
\setbeamercolor{block body example}{fg=primary, bg=success!5}

% ---- Pacotes ----
\usepackage[utf8]{inputenc}
\usepackage[T1]{fontenc}
\usepackage[brazil]{babel}
\usepackage{graphicx}
\usepackage{listings}
\usepackage{xcolor}
\usepackage{tikz}
\usepackage{pifont}
\usepackage{booktabs}
\usepackage{hyperref}

% ---- Ícones simples ----
\newcommand{\cmark}{\textcolor{success}{\ding{51}}}
\newcommand{\xmark}{\textcolor{danger}{\ding{55}}}

% ---- Configuração de código ----
\lstset{
  basicstyle=\ttfamily\scriptsize,
  keywordstyle=\color{primary}\bfseries,
  commentstyle=\color{gray}\itshape,
  stringstyle=\color{accent},
  backgroundcolor=\color{codebg},
  frame=single,
  framesep=3pt,
  rulecolor=\color{primary!15},
  breaklines=true,
  showstringspaces=false,
  tabsize=2,
  numbers=left,
  numberstyle=\tiny\color{gray},
  xleftmargin=1.8em,
  aboveskip=0.4em,
  belowskip=0.2em,
}

% ---- Metadados ----
\title{Ligador, Bibliotecas e Ligação Estática}
\subtitle{Seminário de Software Básico}
\author{Alberto Henrique \quad Caio Veloso \quad Gabriel Ramos \quad Ícaro Vinícius \quad Vinícius Macedo}
\institute{Grupo 2}
\date{2026}

% ============================================================
\begin{document}

% ---- Slide 1: Capa ----
\maketitle

% ---- Slide 2: Roteiro ----
\begin{frame}{Roteiro}
  \tableofcontents
\end{frame}

% ============================================================
\section{Contexto}
% ============================================================

\begin{frame}{Cada Ferramenta Produz a Entrada da Próxima}
  \begin{center}
    \resizebox{0.98\linewidth}{!}{
    \begin{tikzpicture}[
      block/.style={draw=primary!40, fill=lightbg, rounded corners=3pt,
                    minimum width=2.0cm, minimum height=0.75cm,
                    align=center, font=\scriptsize},
      arrow/.style={->, >=stealth, thick, color=primary!60}
    ]
      \node[block] (csrc) at (0,0.8) {Código C\\\texttt{main.c}};
      \node[block] (compiler) at (2.5,0.8) {Compilador\\\texttt{gcc -S}};
      \node[block] (generatedasm) at (5.0,0.8) {Assembly gerado\\\texttt{main.s}};
      \node[block] (asmsrc) at (5.0,-1.2) {Assembly escrito\\\texttt{funcoes.s}};
      \node[block] (assembler) at (7.5,0.8) {Montador\\\texttt{as}};
      \node[block] (objs) at (10.0,0.8) {Arquivos objeto\\\texttt{.o}};
      \node[block, fill=accent!10, draw=accent] (linker) at (12.5,0.8) {\textbf{Ligador}\\(Linker)};
      \node[block] (exec) at (15.0,0.8) {Executável\\ELF};
      \node[block] (loader) at (15.0,-1.2) {Loader\\Processo em memória};
      \node[block] (libs) at (12.5,-1.2) {Bibliotecas\\\texttt{.a} / \texttt{.so}};

      \draw[arrow] (csrc) -- (compiler);
      \draw[arrow] (compiler) -- (generatedasm);
      \draw[arrow] (generatedasm) -- (assembler);
      \draw[arrow] (assembler) -- (objs);
      \draw[arrow] (asmsrc) -- (assembler);
      \draw[arrow] (objs) -- (linker);
      \draw[arrow] (libs) -- (linker);
      \draw[arrow] (linker) -- (exec);
      \draw[arrow] (exec) -- (loader);
    \end{tikzpicture}}
  \end{center}
  \vspace{0.15cm}
  \begin{block}{As fronteiras importam}
    \small O \texttt{gcc} é um driver: com \texttt{-S}, ele para antes do montador.
    O \textbf{linker constrói} o executável; o \textbf{loader} mapeia seus segmentos
    na memória e inicia o processo.
  \end{block}
\end{frame}

% ============================================================
\section{Definições}
% ============================================================

\begin{frame}{O que é o Ligador (Linker)?}
  \begin{columns}[T,onlytextwidth]
    \begin{column}{0.48\textwidth}
      \begin{block}{Definição}
        Programa que \textbf{combina} múltiplos arquivos-objeto e bibliotecas em um \textbf{único executável}.
      \end{block}
      \vspace{0.4cm}
      \begin{description}[\textbf{Saída:}]
        \item[\textbf{Entrada:}] arquivos \texttt{.o} + bibliotecas
        \item[\textbf{Processo:}] resolução de símbolos + relocação
        \item[\textbf{Saída:}] executável binário
      \end{description}
      \vspace{0.2cm}
      {\small\color{gray} Exemplos: \texttt{ld} (GNU), \texttt{link.exe} (MSVC)}
    \end{column}
    \begin{column}{0.48\textwidth}
      \begin{center}
        \resizebox{0.85\linewidth}{!}{
          \begin{tikzpicture}[
            node distance=1.1cm,
            box/.style={draw=primary!30, fill=lightbg, rounded corners=3pt, minimum width=2.1cm, minimum height=0.55cm, align=center, font=\scriptsize},
            linkerbox/.style={draw=accent, fill=accent!10, rounded corners=4pt, minimum width=2.7cm, minimum height=0.75cm, align=center, font=\small\bfseries, text=accent},
            arrow/.style={->, >=stealth, thick, color=primary!50}
          ]
            \node[box] (o1) {main.o};
            \node[box, below of=o1, node distance=0.75cm] (o2) {math.o};
            \node[box, below of=o2, node distance=0.75cm] (a1) {libmath.a};

            \node[linkerbox, right of=o2, node distance=2.7cm] (linker) {Ligador\\(Linker)};

            \node[box, right of=linker, node distance=2.7cm, fill=primary!10, draw=primary] (exec) {programa\\(Executável)};

            \draw[arrow] (o1.east) -- (linker.160);
            \draw[arrow] (o2.east) -- (linker.west);
            \draw[arrow] (a1.east) -- (linker.200);
            \draw[arrow] (linker.east) -- (exec.west);
          \end{tikzpicture}
        }
      \end{center}
    \end{column}
  \end{columns}
\end{frame}

\begin{frame}{Bibliotecas}
  \begin{block}{Definição}
    Coleções de \textbf{código pré-compilado} empacotadas para \textbf{reutilização}.
  \end{block}
  \vspace{0.3cm}
  \begin{columns}[T]
    \begin{column}{0.47\textwidth}
      \begin{exampleblock}{Biblioteca Estática}
        \begin{itemize}
          \item Extensão: \texttt{.a} / \texttt{.lib}
          \item Módulos \textbf{necessários} incorporados ao executável
          \item Criada com \texttt{ar rcs}
        \end{itemize}
      \end{exampleblock}
    \end{column}
    \begin{column}{0.47\textwidth}
      \begin{alertblock}{Biblioteca Dinâmica}
        \begin{itemize}
          \item Extensão: \texttt{.so} / \texttt{.dll}
          \item Carregada em \textbf{runtime}
          \item \textbf{Compartilhada} entre programas
        \end{itemize}
      \end{alertblock}
    \end{column}
  \end{columns}
\end{frame}

\begin{frame}{Ligação Estática}
  \begin{columns}[T]
    \begin{column}{0.53\textwidth}
      \begin{block}{Definição}
        \small Processo realizado em \textbf{tempo de ligação} no qual os
        \textbf{módulos necessários} de uma biblioteca estática são incorporados ao executável.
      \end{block}
      \vspace{-0.05cm}
      {\small
      \begin{itemize}
        \mreducelistspacing
        \item A biblioteca \texttt{.a} não precisa acompanhar o programa.
        \item O linker seleciona os objetos que resolvem símbolos pendentes.
        \item O resultado tende a ser mais independente, porém maior.
      \end{itemize}}
      \vspace{-0.08cm}
      {\scriptsize\textbf{Atenção:} usar uma biblioteca \texttt{.a} não torna todo
      o executável estático; outras bibliotecas podem permanecer dinâmicas.}
    \end{column}
    \begin{column}{0.44\textwidth}
      \begin{center}
        \begin{tikzpicture}[
          box/.style={draw=primary!30, fill=lightbg, rounded corners=3pt,
                      minimum width=3.9cm, minimum height=1.55cm,
                      align=center, font=\scriptsize},
          subbox/.style={draw=accent, fill=accent!10, rounded corners=2pt,
                         minimum width=3.25cm, minimum height=0.45cm,
                         align=center, font=\tiny\bfseries}
        ]
          \node[box, label={[font=\scriptsize\bfseries]above:Ligação Estática}] (static) {};
          \node[subbox] at (static.center) [yshift=0.3cm] {Código do Programa};
          \node[subbox] at (static.center) [yshift=-0.3cm] {Objetos necessários da \texttt{.a}};

          \node[box, below of=static, node distance=2.3cm,
                label={[font=\scriptsize\bfseries]above:Ligação Dinâmica}] (dynamic) {};
          \node[subbox] at (dynamic.center) [yshift=0.3cm] {Código do Programa};
          \node[subbox, fill=white, draw=gray, dashed] at (dynamic.center) [yshift=-0.3cm] {Referência a Biblioteca \texttt{.so}};
        \end{tikzpicture}
      \end{center}
    \end{column}
  \end{columns}
\end{frame}

% ============================================================
\section{Processos e Etapas}
% ============================================================

\begin{frame}{Etapas do Ligador}
  \vspace{-0.2cm}
  \begin{center}
    \begin{tikzpicture}[
      stepbox/.style={draw=primary!30, fill=lightbg, rounded corners=4pt,
                   minimum width=3.2cm, minimum height=0.9cm,
                   align=center, font=\small\bfseries, text=primary},
      num/.style={circle, fill=accent, text=white, font=\bfseries\footnotesize,
                  minimum size=0.55cm, inner sep=0pt},
      arr/.style={->, thick, color=accent!50, >=stealth, shorten >=2pt, shorten <=2pt}
    ]
      \node[num] (n1) at (0, 1) {1};
      \node[stepbox] (s1) at (0,0) {Resolução\\de Símbolos};
      \node[num] (n2) at (4.6, 1) {2};
      \node[stepbox] (s2) at (4.6,0) {Relocação\\de Endereços};
      \node[num] (n3) at (9.2, 1) {3};
      \node[stepbox] (s3) at (9.2,0) {Geração do\\Executável};
      \draw[arr] (s1) -- (s2);
      \draw[arr] (s2) -- (s3);
    \end{tikzpicture}
  \end{center}
  \vspace{0.3cm}
  \begin{enumerate}
    \item \textbf{Resolução de Símbolos}: associa cada referência à sua definição
    \item \textbf{Relocação}: organiza seções e corrige referências dependentes da posição
    \item \textbf{Geração}: produz o ELF final, seus segmentos e o ponto de entrada
  \end{enumerate}
\end{frame}

\begin{frame}{Resolução de Símbolos}
  \begin{columns}[T]
    \begin{column}{0.46\textwidth}
      \begin{block}{Como funciona?}
        \scriptsize
        \begin{itemize}
          \item Cada \texttt{.o} contém uma \textbf{tabela de símbolos}.
          \item Um \textbf{global definido} pode resolver referências externas;
                um símbolo local só vale no próprio módulo.
          \item Um \textbf{indefinido} ainda precisa de uma definição externa.
        \end{itemize}
      \end{block}
      \vspace{0.05cm}
      \begin{alertblock}{Quando a resolução falha}
        \scriptsize
        \texttt{undefined reference}: nenhuma definição encontrada.\\
        \texttt{multiple definition}: mais de uma definição global forte.\\
      \end{alertblock}
    \end{column}
    \begin{column}{0.51\textwidth}
      \begin{center}
        \begin{tikzpicture}[
          filebox/.style={draw=primary!30, fill=lightbg, rounded corners=3pt, minimum width=3.5cm, minimum height=1.6cm, align=center, font=\scriptsize},
          symbol/.style={font=\ttfamily\tiny, align=left}
        ]
          \node[filebox, label={[font=\scriptsize\bfseries]above:main.o}] (main) {};
          \node[symbol] at (main.center) {
            Chamada: \textbf{somar()}\\
            \textcolor{red}{Símbolo Indefinido}
          };

          \node[filebox, below of=main, node distance=2.2cm, label={[font=\scriptsize\bfseries]above:math\_utils.o}] (math) {};
          \node[symbol] at (math.center) {
            Definição: \textbf{somar()}\\
            \textcolor{green!50!black}{Símbolo Definido}
          };

          \draw[->, >=stealth, ultra thick, color=accent, out=180, in=180, looseness=1.4] (main.west) to node[left, font=\scriptsize\bfseries, text=accent] {Resolução} (math.west);
        \end{tikzpicture}
      \end{center}
    \end{column}
  \end{columns}
\end{frame}

\begin{frame}{Relocação de Endereços}
  \begin{block}{Problema}
    Um arquivo \texttt{.o} é \textbf{relocável}: contém offsets relativos às suas seções,
    símbolos e registros que indicam quais posições deverão ser corrigidas pelo linker.
  \end{block}
  \vspace{0.25cm}
  \begin{columns}[T]
    \begin{column}{0.47\textwidth}
      \begin{exampleblock}{Antes \normalfont(referências relocáveis)}
        \texttt{main.o}: \texttt{call somar} $\rightarrow$ destino pendente\\
        \texttt{math.o}: \texttt{somar} no offset \texttt{0x0000} de sua \texttt{.text}
      \end{exampleblock}
    \end{column}
    \begin{column}{0.47\textwidth}
      \begin{block}{Depois \normalfont(layout final)}
        \texttt{main.o}: \texttt{call} com deslocamento corrigido\\
        \texttt{math.o}: função posicionada, por exemplo, em \texttt{0x401020}
      \end{block}
    \end{column}
  \end{columns}
  \vspace{0.3cm}
  \begin{center}
    O linker \alert{organiza as seções}, atribui posições finais e aplica as relocações.
    Em AMD64, muitas referências finais são relativas ao registrador de instruções (RIP).
  \end{center}
\end{frame}

\begin{frame}{Geração do Executável}
  \begin{columns}[T]
    \begin{column}{0.55\textwidth}
      \begin{block}{Arquivo gerado}
        Formato \textbf{ELF} (Linux) ou \textbf{PE} (Windows):
      \end{block}
      \vspace{0.15cm}
      {\small
      \begin{itemize}
        \item \textbf{ELF Header}: metadados e ponto de entrada
        \item \textbf{Program Headers}: segmentos que o loader mapeia na memória
        \item \textbf{\texttt{.text}} e \textbf{\texttt{.rodata}}: código e constantes
        \item \textbf{\texttt{.data}} e \textbf{\texttt{.bss}}: dados inicializados e não inicializados
        \item \textbf{Tabela de símbolos}: útil para análise, mas pode ser removida com \texttt{strip}
      \end{itemize}
      }
    \end{column}
    \begin{column}{0.40\textwidth}
      \begin{center}
      \begin{tikzpicture}[
        sec/.style={draw=primary!25, fill=lightbg, minimum width=4.0cm,
                    minimum height=0.58cm, align=center, font=\scriptsize}
      ]
        \node[sec, fill=primary, font=\scriptsize\color{white}\bfseries] at (0,3.45) {ELF Header};
        \node[sec, fill=accent!10] at (0,2.75) {Program Headers};
        \node[sec] at (0,2.05) {\texttt{.text} / \texttt{.rodata}};
        \node[sec] at (0,1.35) {\texttt{.data} / \texttt{.bss}};
        \node[sec, fill=gray!10] at (0,0.65) {Tabela de símbolos (opcional)};
      \end{tikzpicture}
      \end{center}
    \end{column}
  \end{columns}
\end{frame}

% ============================================================
\section{Vantagens e Desvantagens}
% ============================================================

\begin{frame}{Ligação Estática - Prós e Contras}
  \begin{columns}[T]
    \begin{column}{0.47\textwidth}
      \begin{exampleblock}{Vantagens}
        \begin{itemize}
          \item \textbf{Distribuição simplificada} em sistemas compatíveis
          \item Menor risco de biblioteca \textbf{ausente ou incompatível}
          \item \textbf{Versão controlada} da biblioteca incorporada
          \item Inicialização potencialmente mais simples, sem parte da resolução dinâmica
        \end{itemize}
      \end{exampleblock}
    \end{column}
    \begin{column}{0.47\textwidth}
      \begin{alertblock}{Desvantagens}
        \begin{itemize}
          \item O executável \textbf{tende a ser maior} que um equivalente com bibliotecas compartilhadas
          \item Pode haver \textbf{duplicação} de código entre programas
          \item Atualizar código incorporado exige \textbf{religação e redistribuição}
          \item Algumas bibliotecas de sistema podem não oferecer versão estática
        \end{itemize}
      \end{alertblock}
    \end{column}
  \end{columns}
  \vspace{0.10cm}
  \begin{center}
    {\scriptsize\color{gray}
    Casos comuns: sistemas embarcados e utilitários de linha de comando.
    Containers e binários Go/Rust dependem da configuração e das bibliotecas utilizadas.}
  \end{center}
\end{frame}

% ============================================================
\section{Exemplo Prático}
% ============================================================

\begin{frame}[fragile]{Código Fonte: C (Interface)}
  \begin{columns}[T]
    \begin{column}{0.48\textwidth}
      \begin{block}{\texttt{string\_utils.h}}
        \begin{lstlisting}[language=C]
#ifndef STRING_UTILS_H
#define STRING_UTILS_H

long long my_strlen(const char *str);
void my_reverse(char *str);

#endif
        \end{lstlisting}
      \end{block}
    \end{column}
    \begin{column}{0.48\textwidth}
      \begin{block}{\texttt{main.c}}
        \begin{lstlisting}[language=C]
#include <stdio.h>
#include "string_utils.h"

int main() {
    char texto[] = "Software Basico";
    long long len = my_strlen(texto);
    printf("Tamanho: %lld\n", len);
    my_reverse(texto);
    printf("Invertido: %s\n", texto);
    return 0;
}
        \end{lstlisting}
      \end{block}
    \end{column}
  \end{columns}
\end{frame}

\begin{frame}[fragile]{A ABI Conecta o Código C ao Assembly}
  \begin{columns}[T]
    \begin{column}{0.52\textwidth}
      \begin{block}{\texttt{my\_strlen.s}}
        \begin{lstlisting}[language=C,basicstyle=\ttfamily\scriptsize]
.section .text
.globl my_strlen
.type my_strlen, @function
my_strlen:
    movq $0, %rax
.loop:
    cmpb $0, (%rdi,%rax)
    je .end
    incq %rax
    jmp .loop
.end:
    ret
        \end{lstlisting}
      \end{block}
    \end{column}
    \begin{column}{0.44\textwidth}
      \begin{exampleblock}{ABI System V AMD64}
        \begin{itemize}\small
          \item \texttt{\%rdi}: primeiro argumento, o ponteiro para a string
          \item \texttt{\%rax}: valor retornado por \texttt{my\_strlen}
          \item \texttt{\%rax}, \texttt{\%rcx}, \texttt{\%rdx},
                \texttt{\%rsi} e \texttt{\%rdi} são caller-saved
        \end{itemize}
      \end{exampleblock}
      \vspace{0.2cm}
      {\small \texttt{my\_reverse} também recebe a string em \texttt{\%rdi},
      altera o buffer diretamente e não possui valor de retorno.}
    \end{column}
  \end{columns}
\end{frame}

\begin{frame}[fragile]{\texttt{my\_reverse} Busca as Extremidades e Troca os Bytes}
  \begin{columns}[T]
    \begin{column}{0.48\textwidth}
      \begin{block}{1. Localizar o fim}
        \begin{lstlisting}[language=C,basicstyle=\ttfamily\scriptsize]
.section .text
.globl my_reverse
.type my_reverse, @function
my_reverse:
    movq %rdi, %rsi
.find_end:
    cmpb $0, (%rsi)
    je .prepare
    incq %rsi
    jmp .find_end
.prepare:
    cmpq %rdi, %rsi
    je .done
    decq %rsi
    movq %rdi, %rdx
        \end{lstlisting}
      \end{block}
    \end{column}
    \begin{column}{0.48\textwidth}
      \begin{block}{2. Trocar os bytes}
        \begin{lstlisting}[language=C,basicstyle=\ttfamily\scriptsize]
.swap:
    cmpq %rsi, %rdx
    jge .done
    movb (%rdx), %al
    movb (%rsi), %cl
    movb %cl, (%rdx)
    movb %al, (%rsi)
    incq %rdx
    decq %rsi
    jmp .swap
.done:
    ret
        \end{lstlisting}
      \end{block}
      \vspace{0.05cm}
      {\scriptsize O código usa registradores caller-saved e testa a string
      vazia antes de \texttt{decq}.}
    \end{column}
  \end{columns}
\end{frame}

\begin{frame}[fragile]{Processo de Ligação Estática}
  \begin{block}{1. Montagem e criação da biblioteca \texttt{.a}}
    \begin{lstlisting}[language=bash]
$ as my_strlen.s -o my_strlen.o
$ as my_reverse.s -o my_reverse.o
# nao_usada.s exporta funcao_nao_usada e executa apenas ret
$ as nao_usada.s -o nao_usada.o
$ ar rcs libstrutils.a my_strlen.o my_reverse.o nao_usada.o
    \end{lstlisting}
  \end{block}

  \begin{block}{2. Programa C e duas formas de ligação}
    \begin{lstlisting}[language=bash]
$ gcc -c main.c -o main.o
# somente libstrutils.a ligada estaticamente; libc pode ser dinamica
$ gcc main.o ./libstrutils.a -o programa_misto
# executavel completamente estatico (requer libc estatica instalada)
$ gcc -static main.o ./libstrutils.a -o programa
    \end{lstlisting}
  \end{block}

  \vspace{0.10cm}
  \begin{center}
    {\small A ordem importa: o objeto que cria as referências aparece antes da biblioteca que as resolve.}
  \end{center}
\end{frame}

\begin{frame}[fragile]{Saída e Verificação}
  \begin{columns}[T]
    \begin{column}{0.48\textwidth}
      \begin{block}{Antes da ligação}
        \begin{lstlisting}[language=bash,basicstyle=\ttfamily\tiny]
$ nm main.o | grep -E "my_strlen|my_reverse"
                 U my_reverse
                 U my_strlen
        \end{lstlisting}
      \end{block}
      \begin{block}{Execução}
        \begin{lstlisting}[language=bash,basicstyle=\ttfamily\tiny]
$ ./programa
Tamanho: 15
Invertido: ocisaB erawtfoS
        \end{lstlisting}
      \end{block}
      \begin{block}{Tipo do arquivo}
        \begin{lstlisting}[language=bash,basicstyle=\ttfamily\tiny]
$ file programa
ELF 64-bit LSB executable, x86-64,
statically linked
        \end{lstlisting}
      \end{block}
    \end{column}
    \begin{column}{0.48\textwidth}
      \begin{block}{Depois da ligação - endereços variam}
        \begin{lstlisting}[language=bash,basicstyle=\ttfamily\tiny]
$ nm programa | grep -E "my_strlen|my_reverse"
00000000004010a0 T my_strlen
00000000004010c0 T my_reverse
        \end{lstlisting}
      \end{block}
      \begin{block}{Módulo não utilizado}
        \begin{lstlisting}[language=bash,basicstyle=\ttfamily\tiny]
$ nm programa | grep funcao_nao_usada
# nenhuma saida
        \end{lstlisting}
      \end{block}
      \begin{block}{Verificação com \texttt{ldd}}
        \begin{lstlisting}[language=bash,basicstyle=\ttfamily\tiny]
$ ldd programa
    not a dynamic executable
        \end{lstlisting}
      \end{block}
    \end{column}
  \end{columns}
\end{frame}

\begin{frame}[fragile]{Dentro da Biblioteca \texttt{.a}}
  \begin{columns}[T]
    \begin{column}{0.49\textwidth}
      \begin{block}{Conteúdo do arquivo}
        \begin{lstlisting}[language=bash,basicstyle=\ttfamily\tiny]
$ ar t libstrutils.a
my_strlen.o
my_reverse.o
nao_usada.o
        \end{lstlisting}
      \end{block}
      \begin{block}{Resumo dos símbolos exportados}
        \begin{lstlisting}[language=bash,basicstyle=\ttfamily\tiny]
$ nm -g --defined-only libstrutils.a
my_strlen.o:
0000000000000000 T my_strlen
my_reverse.o:
0000000000000000 T my_reverse
nao_usada.o:
0000000000000000 T funcao_nao_usada
        \end{lstlisting}
      \end{block}
    \end{column}
    \begin{column}{0.47\textwidth}
      \begin{center}
        \vspace{0.15cm}
        \begin{tikzpicture}[
          box/.style={draw=primary!25, fill=lightbg, rounded corners=3pt,
                      minimum width=3.8cm, minimum height=0.60cm,
                      align=center, font=\scriptsize}
        ]
          \node[box, fill=primary, font=\scriptsize\color{white}\bfseries] (lib) at (0,3.1) {libstrutils.a};
          \node[box] (o1) at (0,2.1) {my\_strlen.o};
          \node[box] (o2) at (0,1.3) {my\_reverse.o};
          \node[box, fill=danger!7, draw=danger!45] (o3) at (0,0.5) {nao\_usada.o};
          \draw[->, thick, accent] (lib) -- (o1);
          \draw[->, thick, accent] (lib) -- (o2);
          \draw[->, thick, accent] (lib) -- (o3);
        \end{tikzpicture}
      \end{center}
      \vspace{0.15cm}
      \begin{exampleblock}{Seleção por módulo}
        Como nenhum símbolo de \texttt{nao\_usada.o} é referenciado,
        o linker não extrai esse módulo da biblioteca.
      \end{exampleblock}
    \end{column}
  \end{columns}
\end{frame}

% ============================================================
\section{Considerações Finais}
% ============================================================

\begin{frame}{Resumo}
  \begin{center}
    \begin{tikzpicture}[
      card/.style={draw=primary!20, fill=lightbg, rounded corners=6pt,
                   minimum width=3.4cm, minimum height=2.2cm,
                   align=center, font=\small, text=primary},
    ]
      \node[card] at (0,0) {\textbf{\large Ligador}\\[0.3cm]\small Combina \texttt{.o}'s\\em executável};
      \node[card] at (4.8,0) {\textbf{\large Biblioteca}\\[0.3cm]\small Código pré-compilado\\para reutilização};
      \node[card] at (9.6,0) {\textbf{\large Lig. Estática}\\[0.3cm]\small Incorpora módulos\\necessários da \texttt{.a}};
    \end{tikzpicture}
  \end{center}
  \vspace{0.5cm}
  \begin{itemize}
    \item O Linker é \textbf{essencial} no processo de construção do programa
    \item Ligação estática $\rightarrow$ distribuição mais simples em sistemas \textbf{compatíveis}
    \item Trade-off frequente: possível aumento de \textbf{tamanho} vs menor dependência em tempo de execução
  \end{itemize}
\end{frame}

\begin{frame}{Referências}
  \begin{itemize}\small
    \item STALLINGS, William. \textit{Arquitetura e Organização de Computadores}. 8ª ed. Pearson, 2010.
    \item Material da disciplina de Software Básico: representação de dados, Assembly, pilha, procedimentos e chamadas de sistema.
    \item BRYANT, R. E.; O'HALLARON, D. R. \textit{Computer Systems: A Programmer's Perspective}. 3ª ed. Pearson, 2015. Cap. 7.
    \item LEVINE, J. R. \textit{Linkers and Loaders}. Morgan Kaufmann, 1999.
    \item GNU Binutils: documentação de \texttt{ld}, \texttt{ar}, \texttt{nm} e \texttt{readelf}. \url{https://www.gnu.org/software/binutils/}
    \item \textit{System V Application Binary Interface: AMD64 Architecture Processor Supplement}.
  \end{itemize}
\end{frame}

{
\setbeamercolor{background canvas}{bg=primary}
\begin{frame}[plain]
  \begin{center}
    \vspace{2.5cm}
    {\Huge\color{white}\textbf{Perguntas?}}\\[1cm]
    {\large\color{accent!60!white} Obrigado pela atenção!}\\[0.5cm]
  \end{center}
\end{frame}
}

\end{document}
